% Generated from locale/tr/content/history/biographies/alan-turing.tex; do not hand-edit.


\olfileid[tr]{his}{bio}{tur} 

\olsection{Alan Turing}

Alan Turing 23 Haziran 1912'de Londra'nın Maida Vale semtinde doğdu. Kuramsal
bilgisayar biliminin babası sayılır. Turing'in fen bilimlerine ve matematiğe
ilgisi küçük yaşta başladı. Ancak çocukluk yıllarındaki okullarında edebiyata
ve klasiklere ağırlık verildiğinden ilgi alanları iyi karşılık bulmadı. Bu
nedenle okulda başarısız oldu ve birçok öğretmeninden azar işitti.

\olphoto{turing-alan}{Alan Turing}

Turing lisans öğrencisi olarak Cambridge'deki King's College'a gitti ve
matematik okudu. 1936'da, hesaplanabilir fonksiyon kavramını kesin biçimde
tanımlama ve karar probleminin karar verilemezliğini ispatlama girişimi olarak
(bugün Turing makinesi denen) makineyi geliştirdi. Sonuca, kendi lambda
hesabıyla sonucu ispatlayan Alonzo Church ondan önce ulaştı. Yine de Turing'in
makalesi Church'ün sonucuna atıfla yayımlandı. Church, Turing'i Princeton'a
davet etti; Turing 1936--1938 yıllarını burada geçirdi ve doktorasını Church'ün
danışmanlığında aldı.

Mantığa duyduğu ilgiye karşın Turing'in fen bilimlerine yönelik önceki ilgisi
gücünü korudu. Pratik becerileri, II. Dünya Savaşı sırasında Bletchley
Park'taki Britanya kriptanaliz biriminde verdiği hizmette kullanıldı. Turing,
Alman donanma haberleşmesinde kullanılan şifrenin---Enigma kodunun---
çözülmesinde başat bir rol oynadı. İstatistik ve kriptografi uzmanlığı,
elektronik makinelerin kullanılmaya başlamasıyla birlikte, ekibin ``bombe''
adlı bir şifre çözme makinesi yaratarak kodu kırmasını sağladı. Fikirleri,
dünyanın ilk programlanabilir elektronik bilgisayarı Colossus'un yapımına da
yardım etti; Colossus yine Bletchley Park'ta Alman Lorenz şifresini kırmak
için kullanıldı.

Turing eşcinseldi. Buna rağmen 1942'de Bletchley Park'taki takım
arkadaşlarından Joan Clarke'a evlenme teklif etti; fakat daha sonra eşcinsel
olduğunu ona açıklayıp nişanı bozdu. Yaşamı boyunca birkaç sevgilisi oldu;
oysa eşcinsel edimler o sırada Birleşik Krallık'ta suçtu. Turing'in evi
1952'de o sıradaki sevgilisinin bir arkadaşı tarafından soyuldu. Turing polise
şikâyette bulunurken hükûmetin eşcinsel edimleri yasallaştırmak üzere olduğu
izlenimiyle eşcinsel bir ilişkisi olduğunu söyledi. Bu doğru değildi ve ağır
ahlaksızlık suçlamasıyla yargılandı. Turing hapse girmek yerine libidoyu
azaltan bir hormon tedavisini seçti. 8 Haziran 1954'te siyanür doz aşımı
nedeniyle ölü bulundu---büyük olasılıkla intihar etmişti. Kraliçe II.
Elizabeth 2013'te Turing için kraliyet affı çıkardı.

\begin{reading}
Alan Turing'in kapsamlı bir yaşam öyküsü için \citet{Hodges2014} kaynağına
bakın. Turing'in yaşamı ve çalışmaları, 1996'da Turing rolünde Derek Jacobi
ile televizyona uyarlanan \emph{Breaking the Code} oyununa esin kaynağı oldu.
Benedict Cumberbatch ve Keira Knightley'nin oynadığı, Akademi Ödülü'ne aday
\emph{The Imitation Game} filmi de Alan Turing'in yaşamına ve Bletchley
Park'taki yıllarına serbestçe dayanır \citep{Imitation2014}.

\citet{Radiolab2012}, Turing'in yaşamı ve çalışmaları üzerine birkaç podcast
sunar. BBC Horizon'ın \emph{The Strange Life and Death of Dr.~Turing}
belgeseli çevrimiçi izlenebilir \citep{Sykes1992}. \citep{Theelen2012},
Turing'in 2012'deki yüzüncü doğum yılı onuruna yapılmış, çalışan bir LEGO
Turing Makinesi'nin kısa videosudur.

Turing'in Turing makineleri ve karar problemi üzerine özgün makalesi
\citet{Turing1937}'dir.
\end{reading}

